\chapter{Constitutive State Storage Cleanup}

\section{Motivation}

The semantic cleanup of the materials module revealed a persistent mismatch:
the constitutive-site abstraction in \texttt{fall\_n} stored kinematic history
by default for inelastic materials, even though the vast majority of the code
used \texttt{current\_state()} only as the \emph{last committed kinematic
state}.  This had three undesirable consequences:

\begin{itemize}
  \item the default inelastic path allocated and grew storage that was not
        needed by most algorithms;
  \item the API blurred the distinction between ``current committed value'',
        ``full history'', and ``trial state'';
  \item the architecture was poorly prepared for generic constitutive
        integration algorithms requiring trial/accept/reject workflows.
\end{itemize}

This chapter introduces an explicit state-storage vocabulary while preserving
backward compatibility and the static, header-only nature of the hot path.

\section{Storage Roles}

The cleanup distinguishes three different storage roles.

\subsection{CommittedState}

\texttt{CommittedState<S>} stores a single value of type \texttt{S}.  Every
\texttt{update()} overwrites the previously stored state.  This is the correct
default for:

\begin{itemize}
  \item elastic materials;
  \item inelastic constitutive sites whose irreversible memory already lives
        inside the constitutive law;
  \item post-processing paths that only need the last accepted strain-like
        measure at a material site or section site.
\end{itemize}

Architecturally this is just the existing \texttt{ElasticState<S>} under a new
semantic alias.  The implementation remains trivial and fully inlineable.

\subsection{HistoryState}

\texttt{HistoryState<S>} is the explicit opt-in path for storing a full
kinematic history.  It is semantically equivalent to the legacy
\texttt{MemoryState<S>}, but the key clarification is that the policy now owns
a \emph{storage backend}, not merely a hard-coded \texttt{std::vector}.  Each
\texttt{update()} appends to the supplied storage, and
\texttt{current\_value()} returns the last entry.

Two storage backends are now provided explicitly:

\begin{itemize}
  \item \texttt{DynamicHistoryStorage<S>} for unbounded growth;
  \item \texttt{CircularHistoryStorage<S, N>} for a fixed-size sliding window
        known at compile time.
\end{itemize}

This distinction matters because some constitutive models genuinely need
history, but not \emph{all} of it.  For such cases, a circular buffer provides
deterministic memory, cache-friendly locality, and zero heap growth after
construction.

This path remains useful for:

\begin{itemize}
  \item debugging and validation;
  \item offline material diagnostics;
  \item experimental algorithms that truly need access to all previous
        committed kinematic states.
\end{itemize}

The important change is that \texttt{HistoryState} is no longer the
\emph{implicit} inelastic default, but remains available with whichever
storage backend is appropriate for the algorithm.

\subsection{TrialCommittedState}

\texttt{TrialCommittedState<S>} introduces a minimal two-level storage:

\begin{itemize}
  \item a committed state;
  \item an optional staged trial state.
\end{itemize}

The policy exposes:

\begin{itemize}
  \item \texttt{update()} to stage a trial value;
  \item \texttt{commit\_trial()} to accept it;
  \item \texttt{revert\_trial()} to discard it;
  \item \texttt{current\_value()} to see the trial state if present, otherwise
        the committed one.
\end{itemize}

This policy is not yet the production default.  Its role is preparatory: it
provides the storage semantics needed by future generic integration strategies
without forcing those algorithms to be encoded inside each constitutive law.

\section{MaterialState as a Storage Adapter}

\texttt{MaterialState} now depends on the more accurate concept
\texttt{StateStoragePolicyFor}.  The legacy name
\texttt{MemoryPolicyFor} is preserved as an alias so existing code remains
valid during the transition.

Besides \texttt{current\_value()} and \texttt{update()}, \texttt{MaterialState}
now conditionally forwards:

\begin{itemize}
  \item \texttt{commit\_trial()};
  \item \texttt{revert\_trial()};
  \item \texttt{has\_trial\_value()};
  \item \texttt{committed\_value()};
  \item \texttt{trial\_value()};
  \item \texttt{size()} and indexed logical access when the storage backend
        is history-bearing.
\end{itemize}

This conditional forwarding is implemented with \texttt{requires}-based
introspection.  No virtual dispatch is introduced, and no overhead is paid for
policies that do not support staged trial states.

\section{Change in the Default Inelastic Path}

The most important semantic change is in
\texttt{MaterialInstance<Relation, StatePolicy>}:

\begin{itemize}
  \item \texttt{InelasticConstitutiveSite<R>} now defaults to
        \texttt{CommittedState};
  \item explicit history tracking remains available through
        \texttt{HistoryTrackingConstitutiveSite<R>} and
        \texttt{HistoryTrackedInelasticMaterial<R>};
  \item bounded history tracking is available through
        \texttt{CircularHistoryConstitutiveSite<R, N>}, which binds the site to
        a fixed compile-time circular buffer;
  \item \texttt{TrialConstitutiveSite<R>} exposes the staged-storage path for
        future algorithmic work.
\end{itemize}

This is semantically more honest.  In the current architecture, models such as
\texttt{J2PlasticityRelation}, \texttt{KentParkConcrete}, and
\texttt{MenegottoPintoSteel} store their irreversible internal variables inside
the constitutive relation itself.  The site-level kinematic state is therefore
not the true inelastic memory; it is primarily a local snapshot of the last
accepted measure used by post-processing and reconstruction.

\section{Interaction with Constitutive Updates}

The constitutive-site update path was also adjusted so that storage policies
with staged trial semantics behave correctly under the existing API:

\begin{enumerate}
  \item the constitutive law updates its own irreversible variables;
  \item the site storage receives the kinematic state through
        \texttt{state\_.update(k)};
  \item if the storage policy supports \texttt{commit\_trial()}, the staged
        value is immediately promoted to committed state.
\end{enumerate}

This keeps current behaviour unchanged for production algorithms while making
the storage layer ready for a later refactor in which the constitutive integrator
controls trial and commit phases explicitly.

\section{Why This Matters for Reinforced Concrete Across Scales}

For reinforced concrete across continuum, sectional, and frame scales, a clean
separation between \emph{constitutive law}, \emph{state storage}, and
\emph{integration algorithm} is essential:

\begin{itemize}
  \item frame fiber sections require many local constitutive sites, so storage
        waste multiplies quickly;
  \item dynamic nonlinear analyses require robust and explicit commit/revert
        semantics;
  \item large-displacement and finite-strain formulations need a precise notion
        of trial versus committed state;
  \item continuum-scale cracking and fracture will eventually require local and
        nonlocal internal variables, which should not be confused with merely
        storing a history of kinematic inputs.
\end{itemize}

This cleanup does not yet solve those higher-level problems.  It does,
however, remove a semantic obstacle: site-level storage is no longer forced to
pretend that ``inelastic'' automatically means ``store all past kinematics''.

\section{Next Refactor Stage}

With storage semantics clarified, the next architectural step is now better
defined:

\begin{enumerate}
  \item separate \texttt{ConstitutiveLaw} from \texttt{ConstitutiveState};
  \item let the constitutive integrator evolve the state explicitly instead of
        delegating the entire algorithm to the law;
  \item make the future constitutive-state object choose not only its semantic
        role (committed/history/trial) but also its storage backend
        (dynamic, circular, small-buffer, etc.);
  \item reserve non-owning type erasure for borrowed constitutive views rather
        than for owned site storage;
  \item keep the hot path statically typed for homogeneous loops such as fiber
        integration and section reconstruction.
\end{enumerate}

That next stage is the one that will most directly enable robust nonlinear
reinforced concrete, dynamics, and eventually objective crack/fracture
modelling at the continuum scale.
